\chapter{Conclusiones}
\label{chap:conclusions}

\section{Resultados alcanzados}

El proyecto materializa LIA como una plataforma web que conecta la cartera de
licitaciones con un proceso configurable de cualificación. El núcleo integra
cuentas y roles, pipelines, workflows, oportunidades, catálogos reutilizables,
instancias de cualificación y documentos PDF. Los resultados manuales y
generados por IA comparten el mismo contexto y permanecen revisables por el
usuario.

La arquitectura separa los casos de uso de los adaptadores de persistencia,
almacenamiento e IA. Los estados de ejecución, la evidencia documental y los
reintentos explícitos permiten observar el resultado de las operaciones sin
tratar una generación fallida como una respuesta válida. Los filtros de campos
personalizados amplían la cartera sin duplicar sus filtros nativos y conservan
el ámbito de cuenta.

La evaluación de 92 generaciones fundamenta una recomendación por operación:
Flash para resúmenes y Flash-Lite para preguntas de control, campos
personalizados y decisiones de workflow. En la configuración de un único
modelo global, el ejemplo de despliegue mantiene Flash-Lite por decisión del
desarrollador, mientras que los criterios experimentales recomiendan Flash.
Esta distinción no elimina el riesgo observado en resúmenes ni la necesidad
de revisión humana.

\section{Cumplimiento de objetivos}

La evaluación se refiere a los objetivos registrados y a la evidencia de los
capítulos anteriores, sin extender el alcance a funcionalidades futuras.

\begin{description}
  \item[OG-001.] La plataforma reúne documentación, respuestas y
  pasos de análisis en una oportunidad y permite asistir tareas repetitivas
  mediante IA. Se obtiene una solución funcional para apoyar el análisis y la
  toma de requisitos. Su reducción efectiva de esfuerzo en equipos reales no
  se cuantifica, al no existir una comparación temporal con el trabajo manual.

  \item[OE-001.] Las preguntas de control y el contexto
  documental permiten generar una evaluación preliminar con evidencia
  revisable. El banco experimental demuestra esta capacidad en el conjunto
  controlado, pero no sustituye la acreditación de la empresa ni una evaluación
  profesional sobre cualquier licitación. El objetivo se alcanza dentro de ese
  alcance asistido.

  \item[OE-002.] Se implementan adaptadores sustituibles y un
  banco que relaciona calidad, estructura, latencia, tokens y coste. La
  comparación entre dos modelos de Google permite seleccionar una política por
  operación. La optimización demostrada es experimental: no se atribuye a LIA
  una explotación interna de métricas, un enrutamiento automático por operación
  ni un ahorro mensual observado.

  \item[OE-003.] El editor y la ejecución por oportunidad permiten
  definir pasos, decisiones y acciones, consultar el progreso y completar,
  omitir o reintentar el trabajo. El objetivo se materializa en el núcleo
  implementado, con la exclusión explícita de los ejecutores de tarea y correo.
\end{description}

Las pruebas del \Cref{chap:testing} aportan evidencia de los recorridos
prioritarios y de las fronteras de autorización y cuenta. Las limitaciones del
\Cref{chap:future-work} delimitan esa evidencia: no se afirma una certificación
de seguridad, disponibilidad de producción ni fiabilidad absoluta de la IA.

\section{Valoración personal}
% TODO: incorporar la valoración personal confirmada por el autor.

Desde el punto de vista técnico, el trabajo muestra que integrar IA no consiste
únicamente en invocar un modelo. La utilidad del resultado depende de cómo se
selecciona el contexto, se valida la estructura, se conserva la evidencia y se
gestiona el fallo dentro del proceso de negocio. La distinción entre validez
formal y respaldo documental resulta especialmente relevante en el análisis
de licitaciones.

El diseño modular y las pruebas permiten revisar esas responsabilidades de
forma separada. Al mismo tiempo, introducen un coste de mantenimiento en
contratos, conversiones y estados que debe justificarse por la complejidad del
workflow y el aislamiento multiempresa. La evaluación por operación muestra
también por qué una media global o una tarifa inferior no bastan para elegir
un modelo.

El resultado constituye una base funcional para continuar la validación y las
ampliaciones descritas, manteniendo un principio común: automatizar la
obtención de información no implica delegar la decisión final ni ocultar las
limitaciones de la evidencia disponible.
